\documentclass[11pt]{article}
\usepackage[margin=1in]{geometry}
\usepackage{amsmath,amssymb,amsthm}
\title{Disproof of Conjecture 00000003940}
\author{earthking11}
\date{October 6, 2026}
\newtheorem{proposition}{Proposition}
\begin{document}
\maketitle
\section*{The incompatibility}
The conjecture defines a finite-type invariant as a numerical invariant
obtained from a finite truncation of a differential graded algebra (DGA).  It
then asks for two Legendrians whose DGAs are chain-isomorphic while all their
finite-type invariants differ.

Let $\Phi:A\to B$ be a chain-level DGA isomorphism.  Because $\Phi$ preserves
the differential, multiplication, grading, and filtration, for every truncation
level $n$ it induces an isomorphism
\[
 \Phi_{\le n}:A_{\le n}\xrightarrow{\sim}B_{\le n}.
\]
If $I_n$ is a numerical invariant obtained from the $n$th truncation, invariance
under isomorphism gives
\[
 I_n(A)=I_n(B).
\]
This holds for every $n$.

\begin{proposition}
No two chain-isomorphic DGAs can have all truncation invariants different.
\end{proposition}
\begin{proof}
For an arbitrary truncation level $n$, restrict the chain isomorphism to that
truncation.  Since $I_n$ is an invariant, the displayed equality follows.  It
contradicts the assertion $I_n(A)\ne I_n(B)$, already for one $n$ and hence a
fortiori for all $n$.
\end{proof}

The crossing-number condition ``at most $12$'' is irrelevant: the requested
pair cannot exist at any crossing number.  Calling a numerical construction an
``invariant'' precisely commits it to being constant on the relevant
isomorphism classes.

The accompanying Lean project formalizes this general implication and a
concrete list-truncation model.
\end{document}
